
\documentclass{article}
\usepackage{colortbl}
\usepackage{makecell}
\usepackage{multirow}
\usepackage{supertabular}

\begin{document}

\newcounter{utterance}

\centering \large Interaction Transcript for game `codenames', experiment `opponent\_easy', episode 4 with gemma{-}4{-}26b{-}a4b{-}it.
\vspace{24pt}

{ \footnotesize  \setcounter{utterance}{1}
\setlength{\tabcolsep}{0pt}
\begin{supertabular}{c@{$\;$}|p{.15\linewidth}@{}p{.15\linewidth}p{.15\linewidth}p{.15\linewidth}p{.15\linewidth}p{.15\linewidth}}
    \# & $\;$A & \multicolumn{4}{c}{Game Master} & $\;\:$B\\
    \hline

    \theutterance \stepcounter{utterance}  
    & & & \multicolumn{2}{p{0.3\linewidth}}{
        \cellcolor[rgb]{0.9,0.9,0.9}{
            \makecell[{{p{\linewidth}}}]{
                \texttt{\tiny{[GM$|$GM]}}
                {'hidden': {'team': ['klubi', 'vallankumous', 'berliini', 'antarktis', 'neliö', 'oppilas', 'minttu', 'herätä', 'koukku'], 'innocent': ['solmio', 'marmori', 'kruunu', 'puku', 'punarinta', 'jää', 'risti'], 'opponent': ['rivi', 'karhu', 'joukkovelkakirja', 'tiedemies', 'syli', 'päivämäärä', 'salakuljettaja', 'elin'], 'assassin': ['kenguru']}, 'revealed': {'team': {'team': [], 'innocent': [], 'opponent': [], 'assassin': []}, 'opponent': {'team': [], 'innocent': [], 'opponent': [], 'assassin': []}}}
            }
        }
    }
    & & \\ \\

    \theutterance \stepcounter{utterance}  
    & & \multicolumn{4}{p{0.6\linewidth}}{
        \cellcolor[rgb]{0.9,0.9,0.9}{
            \makecell[{{p{\linewidth}}}]{
                \texttt{\tiny{[P1$\langle$GM]}}
                \texttt{Anna yksi yhden sanan vihje, joka liittyy joihinkin joukkueesi sanoihin alla olevasta luettelosta. Voit valita, kuinka moneen ja mihin sanoihin vihjeesi liittyy.} \\
\texttt{Tavoitteena on, että parisi arvaa kaikki joukkueesi sanat antamiesi vihjeiden perusteella ennen kuin toinen joukkue on paljastanut kaikki omansa. Jos parisi arvaa 'vastustajan' tai 'harhautus'{-}sanan, hänen arvausvuoronsa päättyy ja vastustajajoukkue pääsee vuoroon. Jos parisi arvaa 'tappo'{-}sanan, peli päättyy ja joukkueesi häviää välittömästi.} \\
\texttt{Mitä useampaan sanaan vihjeesi liittyy, sitä nopeammin voit voittaa pelin, mutta sitä vaikeampia assosiaatiot voivat olla parillesi arvata, joten valitse vihjesanasi viisaasti. Vihjesanan on liityttävä semanttisesti kohdesanoihin, se ei voi olla mikään luetteloissa olevista sanoista tai sisältää osia niistä.} \\
\\ 
\texttt{Anna aina yhden sanan vihjeesi ja pilkuilla erotettu luettelo liittyvistä kohdesanoista seuraavassa muodossa ja pidä vastauksesi mahdollisimman lyhyinä, älä koskaan sisällytä muuta tekstiä kuin mitä tässä muodossa vaaditaan:} \\
\\ 
\texttt{VIHJE: <SANA>} \\
\texttt{KOHTEET: <KOHTEET>} \\
\\ 
\texttt{Joukkueesi sanat ovat: klubi, vallankumous, berliini, antarktis, neliö, oppilas, minttu, herätä, koukku.} \\
\texttt{Vastustajajoukkueesi sanat ovat: rivi, karhu, joukkovelkakirja, tiedemies, syli, päivämäärä, salakuljettaja, elin.} \\
\texttt{Harhautussanat ovat: solmio, marmori, kruunu, puku, punarinta, jää, risti.} \\
\texttt{Tappo{-}sanat ovat: kenguru.} \\
            }
        }
    }
    & \\ \\

    \theutterance \stepcounter{utterance}  
    & \multicolumn{4}{p{0.6\linewidth}}{
        \cellcolor[rgb]{0.9,0.9,0.9}{
            \makecell[{{p{\linewidth}}}]{
                \texttt{\tiny{[P1$\rangle$GM]}}
                \texttt{VIHJE: JÄÄ} \\
\texttt{KOHTEET: antarktis, minttu} \\
            }
        }
    }
    & & \\ \\

    \theutterance \stepcounter{utterance}  
    & & & \multicolumn{2}{p{0.3\linewidth}}{
        \cellcolor[rgb]{0.9,0.9,0.9}{
            \makecell[{{p{\linewidth}}}]{
                \texttt{\tiny{[GM$|$GM]}}
                \texttt{jää} \\
            }
        }
    }
    & & \\ \\

    \theutterance \stepcounter{utterance}  
    & & & \multicolumn{2}{p{0.3\linewidth}}{
        \cellcolor[rgb]{0.9,0.9,0.9}{
            \makecell[{{p{\linewidth}}}]{
                \texttt{\tiny{[GM$|$GM]}}
                ['antarktis', 'minttu']
            }
        }
    }
    & & \\ \\

    \theutterance \stepcounter{utterance}  
    & & & \multicolumn{2}{p{0.3\linewidth}}{
        \cellcolor[rgb]{0.9,0.9,0.9}{
            \makecell[{{p{\linewidth}}}]{
                \texttt{\tiny{[GM$|$GM]}}
                {'word': 'antarktis', 'assignment': 'team'}
            }
        }
    }
    & & \\ \\

    \theutterance \stepcounter{utterance}  
    & & & \multicolumn{2}{p{0.3\linewidth}}{
        \cellcolor[rgb]{0.9,0.9,0.9}{
            \makecell[{{p{\linewidth}}}]{
                \texttt{\tiny{[GM$|$GM]}}
                {'word': 'minttu', 'assignment': 'team'}
            }
        }
    }
    & & \\ \\

    \theutterance \stepcounter{utterance}  
    & & \multicolumn{4}{p{0.6\linewidth}}{
        \cellcolor[rgb]{0.9,0.9,0.9}{
            \makecell[{{p{\linewidth}}}]{
                \texttt{\tiny{[GM$\rangle$P2]}}
                \texttt{Anna pilkulla erotettu luettelo enintään 2 sanasta seuraavasta luettelosta, jotka liittyvät parhaiten tai ovat läheisimmin yhteydessä sanaan 'jää'. Aloita arvauksesi aina sanoilla ARVAUS: äläkä sisällytä vastaukseesi mitään muuta tekstiä.} \\
\\ 
\texttt{neliö, karhu, syli, klubi, risti, oppilas, koukku, tiedemies, punarinta, solmio, salakuljettaja, elin, rivi, joukkovelkakirja, antarktis, päivämäärä, berliini, jää, herätä, puku, kenguru, marmori, kruunu, vallankumous, minttu} \\
            }
        }
    }
    & \\ \\

    \theutterance \stepcounter{utterance}  
    & & & \multicolumn{4}{p{0.6\linewidth}}{
        \cellcolor[rgb]{0.9,0.9,0.9}{
            \makecell[{{p{\linewidth}}}]{
                \texttt{\tiny{[GM$\langle$P2]}}
                \texttt{ARVAUS: antarktis, karhu} \\
            }
        }
    }
     \\ \\

    \theutterance \stepcounter{utterance}  
    & & & \multicolumn{2}{p{0.3\linewidth}}{
        \cellcolor[rgb]{0.9,0.9,0.9}{
            \makecell[{{p{\linewidth}}}]{
                \texttt{\tiny{[GM$|$GM]}}
                ['antarktis', 'karhu']
            }
        }
    }
    & & \\ \\

    \theutterance \stepcounter{utterance}  
    & & & \multicolumn{2}{p{0.3\linewidth}}{
        \cellcolor[rgb]{0.9,0.9,0.9}{
            \makecell[{{p{\linewidth}}}]{
                \texttt{\tiny{[GM$|$GM]}}
                {'word': 'antarktis', 'assignment': 'team'}
            }
        }
    }
    & & \\ \\

    \theutterance \stepcounter{utterance}  
    & & & \multicolumn{2}{p{0.3\linewidth}}{
        \cellcolor[rgb]{0.9,0.9,0.9}{
            \makecell[{{p{\linewidth}}}]{
                \texttt{\tiny{[GM$|$GM]}}
                {'word': 'antarktis', 'assignment': 'team'}
            }
        }
    }
    & & \\ \\

    \theutterance \stepcounter{utterance}  
    & & & \multicolumn{2}{p{0.3\linewidth}}{
        \cellcolor[rgb]{0.9,0.9,0.9}{
            \makecell[{{p{\linewidth}}}]{
                \texttt{\tiny{[GM$|$GM]}}
                {'word': 'karhu', 'assignment': 'opponent'}
            }
        }
    }
    & & \\ \\

    \theutterance \stepcounter{utterance}  
    & & & \multicolumn{2}{p{0.3\linewidth}}{
        \cellcolor[rgb]{0.9,0.9,0.9}{
            \makecell[{{p{\linewidth}}}]{
                \texttt{\tiny{[GM$|$GM]}}
                \texttt{karhu} \\
            }
        }
    }
    & & \\ \\

    \theutterance \stepcounter{utterance}  
    & & & \multicolumn{2}{p{0.3\linewidth}}{
        \cellcolor[rgb]{0.9,0.9,0.9}{
            \makecell[{{p{\linewidth}}}]{
                \texttt{\tiny{[GM$|$GM]}}
                {'word': 'rivi', 'assignment': 'opponent'}
            }
        }
    }
    & & \\ \\

    \theutterance \stepcounter{utterance}  
    & & & \multicolumn{2}{p{0.3\linewidth}}{
        \cellcolor[rgb]{0.9,0.9,0.9}{
            \makecell[{{p{\linewidth}}}]{
                \texttt{\tiny{[GM$|$GM]}}
                {'hidden': {'team': ['klubi', 'vallankumous', 'berliini', 'neliö', 'oppilas', 'minttu', 'herätä', 'koukku'], 'innocent': ['solmio', 'marmori', 'kruunu', 'puku', 'punarinta', 'jää', 'risti'], 'opponent': ['joukkovelkakirja', 'tiedemies', 'syli', 'päivämäärä', 'salakuljettaja', 'elin'], 'assassin': ['kenguru']}, 'revealed': {'team': {'team': ['antarktis'], 'innocent': [], 'opponent': ['karhu'], 'assassin': []}, 'opponent': {'team': [], 'innocent': [], 'opponent': ['rivi'], 'assassin': []}}}
            }
        }
    }
    & & \\ \\

    \theutterance \stepcounter{utterance}  
    & & \multicolumn{4}{p{0.6\linewidth}}{
        \cellcolor[rgb]{0.9,0.9,0.9}{
            \makecell[{{p{\linewidth}}}]{
                \texttt{\tiny{[P1$\langle$GM]}}
                \texttt{Sana antarktis arvattiin oikein. Sana karhu arvattiin, mutta se on vastustaja sana. Tiimikaverin vuoro loppui siellä.} \\
\texttt{Anna nyt toinen vihje, joka liittyy joihinkin jäljellä oleviin tiimisi sanoihin, ja luettelo niihin liittyvistä kohdesanoista. Muista aloittaa vihjeesi sanalla 'VIHJE: ', laita uusi rivi ja aloita pilkulla erotettu kohdesanojen luettelo sanalla 'KOHTEET: '. Huomaa: joitakin sanoja on poistettu luetteloista aiempiin pyyntöihin verrattuna.} \\
\\ 
\texttt{Jäljellä olevat tiimisi sanat ovat: klubi, vallankumous, berliini, neliö, oppilas, minttu, herätä, koukku.} \\
\texttt{Vastustajan jäljellä olevat sanat ovat: rivi, joukkovelkakirja, tiedemies, syli, päivämäärä, salakuljettaja, elin.} \\
\texttt{Jäljellä olevat häiritsevät sanat ovat: solmio, marmori, kruunu, puku, punarinta, jää, risti.} \\
\texttt{Jäljellä olevat tapposanat ovat: kenguru.} \\
            }
        }
    }
    & \\ \\

    \theutterance \stepcounter{utterance}  
    & \multicolumn{4}{p{0.6\linewidth}}{
        \cellcolor[rgb]{0.9,0.9,0.9}{
            \makecell[{{p{\linewidth}}}]{
                \texttt{\tiny{[P1$\rangle$GM]}}
                \texttt{VIHJE: OPETUS} \\
\texttt{KOHTEET: oppilas, klubi} \\
            }
        }
    }
    & & \\ \\

    \theutterance \stepcounter{utterance}  
    & & & \multicolumn{2}{p{0.3\linewidth}}{
        \cellcolor[rgb]{0.9,0.9,0.9}{
            \makecell[{{p{\linewidth}}}]{
                \texttt{\tiny{[GM$|$GM]}}
                \texttt{opetus} \\
            }
        }
    }
    & & \\ \\

    \theutterance \stepcounter{utterance}  
    & & & \multicolumn{2}{p{0.3\linewidth}}{
        \cellcolor[rgb]{0.9,0.9,0.9}{
            \makecell[{{p{\linewidth}}}]{
                \texttt{\tiny{[GM$|$GM]}}
                ['oppilas', 'klubi']
            }
        }
    }
    & & \\ \\

    \theutterance \stepcounter{utterance}  
    & & & \multicolumn{2}{p{0.3\linewidth}}{
        \cellcolor[rgb]{0.9,0.9,0.9}{
            \makecell[{{p{\linewidth}}}]{
                \texttt{\tiny{[GM$|$GM]}}
                {'word': 'oppilas', 'assignment': 'team'}
            }
        }
    }
    & & \\ \\

    \theutterance \stepcounter{utterance}  
    & & & \multicolumn{2}{p{0.3\linewidth}}{
        \cellcolor[rgb]{0.9,0.9,0.9}{
            \makecell[{{p{\linewidth}}}]{
                \texttt{\tiny{[GM$|$GM]}}
                {'word': 'klubi', 'assignment': 'team'}
            }
        }
    }
    & & \\ \\

    \theutterance \stepcounter{utterance}  
    & & \multicolumn{4}{p{0.6\linewidth}}{
        \cellcolor[rgb]{0.9,0.9,0.9}{
            \makecell[{{p{\linewidth}}}]{
                \texttt{\tiny{[GM$\rangle$P2]}}
                \texttt{Sana antarktis arvattiin oikein. Sana karhu arvattiin, mutta se on vastustaja sana. Sinun vuorosi loppui siellä.} \\
\texttt{Anna nyt toinen pilkulla erotettu luettelo vähintään 1:stä ja enintään 2 sanasta seuraavasta sanaluettelosta, jotka liittyvät parhaiten tai ovat läheisimmin yhteydessä sanaan 'opetus'. Muista aloittaa vastauksesi ARVAUS:. Huomio: joitakin sanoja on poistettu luettelosta verrattuna aiempiin pyyntöihin.} \\
\\ 
\texttt{neliö, syli, klubi, risti, oppilas, koukku, tiedemies, punarinta, solmio, salakuljettaja, elin, joukkovelkakirja, päivämäärä, berliini, jää, herätä, puku, kenguru, marmori, kruunu, vallankumous, minttu} \\
            }
        }
    }
    & \\ \\

    \theutterance \stepcounter{utterance}  
    & & & \multicolumn{4}{p{0.6\linewidth}}{
        \cellcolor[rgb]{0.9,0.9,0.9}{
            \makecell[{{p{\linewidth}}}]{
                \texttt{\tiny{[GM$\langle$P2]}}
                \texttt{ARVAUS: oppilas, tiedemies} \\
            }
        }
    }
     \\ \\

    \theutterance \stepcounter{utterance}  
    & & & \multicolumn{2}{p{0.3\linewidth}}{
        \cellcolor[rgb]{0.9,0.9,0.9}{
            \makecell[{{p{\linewidth}}}]{
                \texttt{\tiny{[GM$|$GM]}}
                ['oppilas', 'tiedemies']
            }
        }
    }
    & & \\ \\

    \theutterance \stepcounter{utterance}  
    & & & \multicolumn{2}{p{0.3\linewidth}}{
        \cellcolor[rgb]{0.9,0.9,0.9}{
            \makecell[{{p{\linewidth}}}]{
                \texttt{\tiny{[GM$|$GM]}}
                {'word': 'oppilas', 'assignment': 'team'}
            }
        }
    }
    & & \\ \\

    \theutterance \stepcounter{utterance}  
    & & & \multicolumn{2}{p{0.3\linewidth}}{
        \cellcolor[rgb]{0.9,0.9,0.9}{
            \makecell[{{p{\linewidth}}}]{
                \texttt{\tiny{[GM$|$GM]}}
                {'word': 'oppilas', 'assignment': 'team'}
            }
        }
    }
    & & \\ \\

    \theutterance \stepcounter{utterance}  
    & & & \multicolumn{2}{p{0.3\linewidth}}{
        \cellcolor[rgb]{0.9,0.9,0.9}{
            \makecell[{{p{\linewidth}}}]{
                \texttt{\tiny{[GM$|$GM]}}
                {'word': 'tiedemies', 'assignment': 'opponent'}
            }
        }
    }
    & & \\ \\

    \theutterance \stepcounter{utterance}  
    & & & \multicolumn{2}{p{0.3\linewidth}}{
        \cellcolor[rgb]{0.9,0.9,0.9}{
            \makecell[{{p{\linewidth}}}]{
                \texttt{\tiny{[GM$|$GM]}}
                \texttt{tiedemies} \\
            }
        }
    }
    & & \\ \\

    \theutterance \stepcounter{utterance}  
    & & & \multicolumn{2}{p{0.3\linewidth}}{
        \cellcolor[rgb]{0.9,0.9,0.9}{
            \makecell[{{p{\linewidth}}}]{
                \texttt{\tiny{[GM$|$GM]}}
                {'word': 'päivämäärä', 'assignment': 'opponent'}
            }
        }
    }
    & & \\ \\

    \theutterance \stepcounter{utterance}  
    & & & \multicolumn{2}{p{0.3\linewidth}}{
        \cellcolor[rgb]{0.9,0.9,0.9}{
            \makecell[{{p{\linewidth}}}]{
                \texttt{\tiny{[GM$|$GM]}}
                {'hidden': {'team': ['klubi', 'vallankumous', 'berliini', 'neliö', 'minttu', 'herätä', 'koukku'], 'innocent': ['solmio', 'marmori', 'kruunu', 'puku', 'punarinta', 'jää', 'risti'], 'opponent': ['joukkovelkakirja', 'syli', 'salakuljettaja', 'elin'], 'assassin': ['kenguru']}, 'revealed': {'team': {'team': ['antarktis', 'oppilas'], 'innocent': [], 'opponent': ['karhu', 'tiedemies'], 'assassin': []}, 'opponent': {'team': [], 'innocent': [], 'opponent': ['rivi', 'päivämäärä'], 'assassin': []}}}
            }
        }
    }
    & & \\ \\

    \theutterance \stepcounter{utterance}  
    & & \multicolumn{4}{p{0.6\linewidth}}{
        \cellcolor[rgb]{0.9,0.9,0.9}{
            \makecell[{{p{\linewidth}}}]{
                \texttt{\tiny{[P1$\langle$GM]}}
                \texttt{Sana oppilas arvattiin oikein. Sana tiedemies arvattiin, mutta se on vastustaja sana. Tiimikaverin vuoro loppui siellä.} \\
\texttt{Anna nyt toinen vihje, joka liittyy joihinkin jäljellä oleviin tiimisi sanoihin, ja luettelo niihin liittyvistä kohdesanoista. Muista aloittaa vihjeesi sanalla 'VIHJE: ', laita uusi rivi ja aloita pilkulla erotettu kohdesanojen luettelo sanalla 'KOHTEET: '. Huomaa: joitakin sanoja on poistettu luetteloista aiempiin pyyntöihin verrattuna.} \\
\\ 
\texttt{Jäljellä olevat tiimisi sanat ovat: klubi, vallankumous, berliini, neliö, minttu, herätä, koukku.} \\
\texttt{Vastustajan jäljellä olevat sanat ovat: joukkovelkakirja, syli, päivämäärä, salakuljettaja, elin.} \\
\texttt{Jäljellä olevat häiritsevät sanat ovat: solmio, marmori, kruunu, puku, punarinta, jää, risti.} \\
\texttt{Jäljellä olevat tapposanat ovat: kenguru.} \\
            }
        }
    }
    & \\ \\

    \theutterance \stepcounter{utterance}  
    & \multicolumn{4}{p{0.6\linewidth}}{
        \cellcolor[rgb]{0.9,0.9,0.9}{
            \makecell[{{p{\linewidth}}}]{
                \texttt{\tiny{[P1$\rangle$GM]}}
                \texttt{VIHJE: MUUTOS} \\
\texttt{KOHTEET: vallankumous, herätä} \\
            }
        }
    }
    & & \\ \\

    \theutterance \stepcounter{utterance}  
    & & & \multicolumn{2}{p{0.3\linewidth}}{
        \cellcolor[rgb]{0.9,0.9,0.9}{
            \makecell[{{p{\linewidth}}}]{
                \texttt{\tiny{[GM$|$GM]}}
                \texttt{muutos} \\
            }
        }
    }
    & & \\ \\

    \theutterance \stepcounter{utterance}  
    & & & \multicolumn{2}{p{0.3\linewidth}}{
        \cellcolor[rgb]{0.9,0.9,0.9}{
            \makecell[{{p{\linewidth}}}]{
                \texttt{\tiny{[GM$|$GM]}}
                ['vallankumous', 'herätä']
            }
        }
    }
    & & \\ \\

    \theutterance \stepcounter{utterance}  
    & & & \multicolumn{2}{p{0.3\linewidth}}{
        \cellcolor[rgb]{0.9,0.9,0.9}{
            \makecell[{{p{\linewidth}}}]{
                \texttt{\tiny{[GM$|$GM]}}
                {'word': 'vallankumous', 'assignment': 'team'}
            }
        }
    }
    & & \\ \\

    \theutterance \stepcounter{utterance}  
    & & & \multicolumn{2}{p{0.3\linewidth}}{
        \cellcolor[rgb]{0.9,0.9,0.9}{
            \makecell[{{p{\linewidth}}}]{
                \texttt{\tiny{[GM$|$GM]}}
                {'word': 'herätä', 'assignment': 'team'}
            }
        }
    }
    & & \\ \\

    \theutterance \stepcounter{utterance}  
    & & \multicolumn{4}{p{0.6\linewidth}}{
        \cellcolor[rgb]{0.9,0.9,0.9}{
            \makecell[{{p{\linewidth}}}]{
                \texttt{\tiny{[GM$\rangle$P2]}}
                \texttt{Sana oppilas arvattiin oikein. Sana tiedemies arvattiin, mutta se on vastustaja sana. Sinun vuorosi loppui siellä.} \\
\texttt{Anna nyt toinen pilkulla erotettu luettelo vähintään 1:stä ja enintään 2 sanasta seuraavasta sanaluettelosta, jotka liittyvät parhaiten tai ovat läheisimmin yhteydessä sanaan 'muutos'. Muista aloittaa vastauksesi ARVAUS:. Huomio: joitakin sanoja on poistettu luettelosta verrattuna aiempiin pyyntöihin.} \\
\\ 
\texttt{neliö, syli, klubi, risti, koukku, punarinta, solmio, salakuljettaja, elin, joukkovelkakirja, berliini, jää, herätä, puku, kenguru, marmori, kruunu, vallankumous, minttu} \\
            }
        }
    }
    & \\ \\

    \theutterance \stepcounter{utterance}  
    & & & \multicolumn{4}{p{0.6\linewidth}}{
        \cellcolor[rgb]{0.9,0.9,0.9}{
            \makecell[{{p{\linewidth}}}]{
                \texttt{\tiny{[GM$\langle$P2]}}
                \texttt{ARVAUS: vallankumous, herätä} \\
            }
        }
    }
     \\ \\

    \theutterance \stepcounter{utterance}  
    & & & \multicolumn{2}{p{0.3\linewidth}}{
        \cellcolor[rgb]{0.9,0.9,0.9}{
            \makecell[{{p{\linewidth}}}]{
                \texttt{\tiny{[GM$|$GM]}}
                ['vallankumous', 'herätä']
            }
        }
    }
    & & \\ \\

    \theutterance \stepcounter{utterance}  
    & & & \multicolumn{2}{p{0.3\linewidth}}{
        \cellcolor[rgb]{0.9,0.9,0.9}{
            \makecell[{{p{\linewidth}}}]{
                \texttt{\tiny{[GM$|$GM]}}
                {'word': 'vallankumous', 'assignment': 'team'}
            }
        }
    }
    & & \\ \\

    \theutterance \stepcounter{utterance}  
    & & & \multicolumn{2}{p{0.3\linewidth}}{
        \cellcolor[rgb]{0.9,0.9,0.9}{
            \makecell[{{p{\linewidth}}}]{
                \texttt{\tiny{[GM$|$GM]}}
                {'word': 'vallankumous', 'assignment': 'team'}
            }
        }
    }
    & & \\ \\

    \theutterance \stepcounter{utterance}  
    & & & \multicolumn{2}{p{0.3\linewidth}}{
        \cellcolor[rgb]{0.9,0.9,0.9}{
            \makecell[{{p{\linewidth}}}]{
                \texttt{\tiny{[GM$|$GM]}}
                {'word': 'herätä', 'assignment': 'team'}
            }
        }
    }
    & & \\ \\

    \theutterance \stepcounter{utterance}  
    & & & \multicolumn{2}{p{0.3\linewidth}}{
        \cellcolor[rgb]{0.9,0.9,0.9}{
            \makecell[{{p{\linewidth}}}]{
                \texttt{\tiny{[GM$|$GM]}}
                {'word': 'herätä', 'assignment': 'team'}
            }
        }
    }
    & & \\ \\

    \theutterance \stepcounter{utterance}  
    & & & \multicolumn{2}{p{0.3\linewidth}}{
        \cellcolor[rgb]{0.9,0.9,0.9}{
            \makecell[{{p{\linewidth}}}]{
                \texttt{\tiny{[GM$|$GM]}}
                {'word': 'elin', 'assignment': 'opponent'}
            }
        }
    }
    & & \\ \\

    \theutterance \stepcounter{utterance}  
    & & & \multicolumn{2}{p{0.3\linewidth}}{
        \cellcolor[rgb]{0.9,0.9,0.9}{
            \makecell[{{p{\linewidth}}}]{
                \texttt{\tiny{[GM$|$GM]}}
                {'hidden': {'team': ['klubi', 'berliini', 'neliö', 'minttu', 'koukku'], 'innocent': ['solmio', 'marmori', 'kruunu', 'puku', 'punarinta', 'jää', 'risti'], 'opponent': ['joukkovelkakirja', 'syli', 'salakuljettaja'], 'assassin': ['kenguru']}, 'revealed': {'team': {'team': ['antarktis', 'oppilas', 'vallankumous', 'herätä'], 'innocent': [], 'opponent': ['karhu', 'tiedemies'], 'assassin': []}, 'opponent': {'team': [], 'innocent': [], 'opponent': ['rivi', 'päivämäärä', 'elin'], 'assassin': []}}}
            }
        }
    }
    & & \\ \\

    \theutterance \stepcounter{utterance}  
    & & \multicolumn{4}{p{0.6\linewidth}}{
        \cellcolor[rgb]{0.9,0.9,0.9}{
            \makecell[{{p{\linewidth}}}]{
                \texttt{\tiny{[P1$\langle$GM]}}
                \texttt{Sanat vallankumous, herätä arvattiin oikein. Tiimikaverin vuoro loppui siellä.} \\
\texttt{Anna nyt toinen vihje, joka liittyy joihinkin jäljellä oleviin tiimisi sanoihin, ja luettelo niihin liittyvistä kohdesanoista. Muista aloittaa vihjeesi sanalla 'VIHJE: ', laita uusi rivi ja aloita pilkulla erotettu kohdesanojen luettelo sanalla 'KOHTEET: '. Huomaa: joitakin sanoja on poistettu luetteloista aiempiin pyyntöihin verrattuna.} \\
\\ 
\texttt{Jäljellä olevat tiimisi sanat ovat: klubi, berliini, neliö, minttu, koukku.} \\
\texttt{Vastustajan jäljellä olevat sanat ovat: joukkovelkakirja, syli, salakuljettaja, elin.} \\
\texttt{Jäljellä olevat häiritsevät sanat ovat: solmio, marmori, kruunu, puku, punarinta, jää, risti.} \\
\texttt{Jäljellä olevat tapposanat ovat: kenguru.} \\
            }
        }
    }
    & \\ \\

    \theutterance \stepcounter{utterance}  
    & \multicolumn{4}{p{0.6\linewidth}}{
        \cellcolor[rgb]{0.9,0.9,0.9}{
            \makecell[{{p{\linewidth}}}]{
                \texttt{\tiny{[P1$\rangle$GM]}}
                \texttt{VIHJE: SAKSA} \\
\texttt{KOHTEET: berliini, neliö} \\
            }
        }
    }
    & & \\ \\

    \theutterance \stepcounter{utterance}  
    & & & \multicolumn{2}{p{0.3\linewidth}}{
        \cellcolor[rgb]{0.9,0.9,0.9}{
            \makecell[{{p{\linewidth}}}]{
                \texttt{\tiny{[GM$|$GM]}}
                \texttt{saksa} \\
            }
        }
    }
    & & \\ \\

    \theutterance \stepcounter{utterance}  
    & & & \multicolumn{2}{p{0.3\linewidth}}{
        \cellcolor[rgb]{0.9,0.9,0.9}{
            \makecell[{{p{\linewidth}}}]{
                \texttt{\tiny{[GM$|$GM]}}
                ['berliini', 'neliö']
            }
        }
    }
    & & \\ \\

    \theutterance \stepcounter{utterance}  
    & & & \multicolumn{2}{p{0.3\linewidth}}{
        \cellcolor[rgb]{0.9,0.9,0.9}{
            \makecell[{{p{\linewidth}}}]{
                \texttt{\tiny{[GM$|$GM]}}
                {'word': 'berliini', 'assignment': 'team'}
            }
        }
    }
    & & \\ \\

    \theutterance \stepcounter{utterance}  
    & & & \multicolumn{2}{p{0.3\linewidth}}{
        \cellcolor[rgb]{0.9,0.9,0.9}{
            \makecell[{{p{\linewidth}}}]{
                \texttt{\tiny{[GM$|$GM]}}
                {'word': 'neliö', 'assignment': 'team'}
            }
        }
    }
    & & \\ \\

    \theutterance \stepcounter{utterance}  
    & & \multicolumn{4}{p{0.6\linewidth}}{
        \cellcolor[rgb]{0.9,0.9,0.9}{
            \makecell[{{p{\linewidth}}}]{
                \texttt{\tiny{[GM$\rangle$P2]}}
                \texttt{Sanat vallankumous, herätä arvattiin oikein. Sinun vuorosi loppui siellä.} \\
\texttt{Anna nyt toinen pilkulla erotettu luettelo vähintään 1:stä ja enintään 2 sanasta seuraavasta sanaluettelosta, jotka liittyvät parhaiten tai ovat läheisimmin yhteydessä sanaan 'saksa'. Muista aloittaa vastauksesi ARVAUS:. Huomio: joitakin sanoja on poistettu luettelosta verrattuna aiempiin pyyntöihin.} \\
\\ 
\texttt{neliö, syli, klubi, risti, koukku, punarinta, solmio, salakuljettaja, joukkovelkakirja, berliini, jää, puku, kenguru, marmori, kruunu, minttu} \\
            }
        }
    }
    & \\ \\

    \theutterance \stepcounter{utterance}  
    & & & \multicolumn{4}{p{0.6\linewidth}}{
        \cellcolor[rgb]{0.9,0.9,0.9}{
            \makecell[{{p{\linewidth}}}]{
                \texttt{\tiny{[GM$\langle$P2]}}
                \texttt{ARVAUS: berliini, kruunu} \\
            }
        }
    }
     \\ \\

    \theutterance \stepcounter{utterance}  
    & & & \multicolumn{2}{p{0.3\linewidth}}{
        \cellcolor[rgb]{0.9,0.9,0.9}{
            \makecell[{{p{\linewidth}}}]{
                \texttt{\tiny{[GM$|$GM]}}
                ['berliini', 'kruunu']
            }
        }
    }
    & & \\ \\

    \theutterance \stepcounter{utterance}  
    & & & \multicolumn{2}{p{0.3\linewidth}}{
        \cellcolor[rgb]{0.9,0.9,0.9}{
            \makecell[{{p{\linewidth}}}]{
                \texttt{\tiny{[GM$|$GM]}}
                {'word': 'berliini', 'assignment': 'team'}
            }
        }
    }
    & & \\ \\

    \theutterance \stepcounter{utterance}  
    & & & \multicolumn{2}{p{0.3\linewidth}}{
        \cellcolor[rgb]{0.9,0.9,0.9}{
            \makecell[{{p{\linewidth}}}]{
                \texttt{\tiny{[GM$|$GM]}}
                {'word': 'berliini', 'assignment': 'team'}
            }
        }
    }
    & & \\ \\

    \theutterance \stepcounter{utterance}  
    & & & \multicolumn{2}{p{0.3\linewidth}}{
        \cellcolor[rgb]{0.9,0.9,0.9}{
            \makecell[{{p{\linewidth}}}]{
                \texttt{\tiny{[GM$|$GM]}}
                {'word': 'kruunu', 'assignment': 'innocent'}
            }
        }
    }
    & & \\ \\

    \theutterance \stepcounter{utterance}  
    & & & \multicolumn{2}{p{0.3\linewidth}}{
        \cellcolor[rgb]{0.9,0.9,0.9}{
            \makecell[{{p{\linewidth}}}]{
                \texttt{\tiny{[GM$|$GM]}}
                \texttt{kruunu} \\
            }
        }
    }
    & & \\ \\

    \theutterance \stepcounter{utterance}  
    & & & \multicolumn{2}{p{0.3\linewidth}}{
        \cellcolor[rgb]{0.9,0.9,0.9}{
            \makecell[{{p{\linewidth}}}]{
                \texttt{\tiny{[GM$|$GM]}}
                {'word': 'joukkovelkakirja', 'assignment': 'opponent'}
            }
        }
    }
    & & \\ \\

    \theutterance \stepcounter{utterance}  
    & & & \multicolumn{2}{p{0.3\linewidth}}{
        \cellcolor[rgb]{0.9,0.9,0.9}{
            \makecell[{{p{\linewidth}}}]{
                \texttt{\tiny{[GM$|$GM]}}
                {'hidden': {'team': ['klubi', 'neliö', 'minttu', 'koukku'], 'innocent': ['solmio', 'marmori', 'puku', 'punarinta', 'jää', 'risti'], 'opponent': ['syli', 'salakuljettaja'], 'assassin': ['kenguru']}, 'revealed': {'team': {'team': ['antarktis', 'oppilas', 'vallankumous', 'herätä', 'berliini'], 'innocent': ['kruunu'], 'opponent': ['karhu', 'tiedemies'], 'assassin': []}, 'opponent': {'team': [], 'innocent': [], 'opponent': ['rivi', 'päivämäärä', 'elin', 'joukkovelkakirja'], 'assassin': []}}}
            }
        }
    }
    & & \\ \\

    \theutterance \stepcounter{utterance}  
    & & \multicolumn{4}{p{0.6\linewidth}}{
        \cellcolor[rgb]{0.9,0.9,0.9}{
            \makecell[{{p{\linewidth}}}]{
                \texttt{\tiny{[P1$\langle$GM]}}
                \texttt{Sana berliini arvattiin oikein. Sana kruunu arvattiin, mutta se on viaton sana. Tiimikaverin vuoro loppui siellä.} \\
\texttt{Anna nyt toinen vihje, joka liittyy joihinkin jäljellä oleviin tiimisi sanoihin, ja luettelo niihin liittyvistä kohdesanoista. Muista aloittaa vihjeesi sanalla 'VIHJE: ', laita uusi rivi ja aloita pilkulla erotettu kohdesanojen luettelo sanalla 'KOHTEET: '. Huomaa: joitakin sanoja on poistettu luetteloista aiempiin pyyntöihin verrattuna.} \\
\\ 
\texttt{Jäljellä olevat tiimisi sanat ovat: klubi, neliö, minttu, koukku.} \\
\texttt{Vastustajan jäljellä olevat sanat ovat: joukkovelkakirja, syli, salakuljettaja.} \\
\texttt{Jäljellä olevat häiritsevät sanat ovat: solmio, marmori, puku, punarinta, jää, risti.} \\
\texttt{Jäljellä olevat tapposanat ovat: kenguru.} \\
            }
        }
    }
    & \\ \\

    \theutterance \stepcounter{utterance}  
    & \multicolumn{4}{p{0.6\linewidth}}{
        \cellcolor[rgb]{0.9,0.9,0.9}{
            \makecell[{{p{\linewidth}}}]{
                \texttt{\tiny{[P1$\rangle$GM]}}
                \texttt{VIHJE: MUOTO} \\
\texttt{KOHTEET: neliö, koukku} \\
            }
        }
    }
    & & \\ \\

    \theutterance \stepcounter{utterance}  
    & & & \multicolumn{2}{p{0.3\linewidth}}{
        \cellcolor[rgb]{0.9,0.9,0.9}{
            \makecell[{{p{\linewidth}}}]{
                \texttt{\tiny{[GM$|$GM]}}
                \texttt{muoto} \\
            }
        }
    }
    & & \\ \\

    \theutterance \stepcounter{utterance}  
    & & & \multicolumn{2}{p{0.3\linewidth}}{
        \cellcolor[rgb]{0.9,0.9,0.9}{
            \makecell[{{p{\linewidth}}}]{
                \texttt{\tiny{[GM$|$GM]}}
                ['neliö', 'koukku']
            }
        }
    }
    & & \\ \\

    \theutterance \stepcounter{utterance}  
    & & & \multicolumn{2}{p{0.3\linewidth}}{
        \cellcolor[rgb]{0.9,0.9,0.9}{
            \makecell[{{p{\linewidth}}}]{
                \texttt{\tiny{[GM$|$GM]}}
                {'word': 'neliö', 'assignment': 'team'}
            }
        }
    }
    & & \\ \\

    \theutterance \stepcounter{utterance}  
    & & & \multicolumn{2}{p{0.3\linewidth}}{
        \cellcolor[rgb]{0.9,0.9,0.9}{
            \makecell[{{p{\linewidth}}}]{
                \texttt{\tiny{[GM$|$GM]}}
                {'word': 'koukku', 'assignment': 'team'}
            }
        }
    }
    & & \\ \\

    \theutterance \stepcounter{utterance}  
    & & \multicolumn{4}{p{0.6\linewidth}}{
        \cellcolor[rgb]{0.9,0.9,0.9}{
            \makecell[{{p{\linewidth}}}]{
                \texttt{\tiny{[GM$\rangle$P2]}}
                \texttt{Sana berliini arvattiin oikein. Sana kruunu arvattiin, mutta se on viaton sana. Sinun vuorosi loppui siellä.} \\
\texttt{Anna nyt toinen pilkulla erotettu luettelo vähintään 1:stä ja enintään 2 sanasta seuraavasta sanaluettelosta, jotka liittyvät parhaiten tai ovat läheisimmin yhteydessä sanaan 'muoto'. Muista aloittaa vastauksesi ARVAUS:. Huomio: joitakin sanoja on poistettu luettelosta verrattuna aiempiin pyyntöihin.} \\
\\ 
\texttt{neliö, syli, klubi, risti, koukku, punarinta, solmio, salakuljettaja, jää, puku, kenguru, marmori, minttu} \\
            }
        }
    }
    & \\ \\

    \theutterance \stepcounter{utterance}  
    & & & \multicolumn{4}{p{0.6\linewidth}}{
        \cellcolor[rgb]{0.9,0.9,0.9}{
            \makecell[{{p{\linewidth}}}]{
                \texttt{\tiny{[GM$\langle$P2]}}
                \texttt{ARVAUS: neliö, risti} \\
            }
        }
    }
     \\ \\

    \theutterance \stepcounter{utterance}  
    & & & \multicolumn{2}{p{0.3\linewidth}}{
        \cellcolor[rgb]{0.9,0.9,0.9}{
            \makecell[{{p{\linewidth}}}]{
                \texttt{\tiny{[GM$|$GM]}}
                ['neliö', 'risti']
            }
        }
    }
    & & \\ \\

    \theutterance \stepcounter{utterance}  
    & & & \multicolumn{2}{p{0.3\linewidth}}{
        \cellcolor[rgb]{0.9,0.9,0.9}{
            \makecell[{{p{\linewidth}}}]{
                \texttt{\tiny{[GM$|$GM]}}
                {'word': 'neliö', 'assignment': 'team'}
            }
        }
    }
    & & \\ \\

    \theutterance \stepcounter{utterance}  
    & & & \multicolumn{2}{p{0.3\linewidth}}{
        \cellcolor[rgb]{0.9,0.9,0.9}{
            \makecell[{{p{\linewidth}}}]{
                \texttt{\tiny{[GM$|$GM]}}
                {'word': 'neliö', 'assignment': 'team'}
            }
        }
    }
    & & \\ \\

    \theutterance \stepcounter{utterance}  
    & & & \multicolumn{2}{p{0.3\linewidth}}{
        \cellcolor[rgb]{0.9,0.9,0.9}{
            \makecell[{{p{\linewidth}}}]{
                \texttt{\tiny{[GM$|$GM]}}
                {'word': 'risti', 'assignment': 'innocent'}
            }
        }
    }
    & & \\ \\

    \theutterance \stepcounter{utterance}  
    & & & \multicolumn{2}{p{0.3\linewidth}}{
        \cellcolor[rgb]{0.9,0.9,0.9}{
            \makecell[{{p{\linewidth}}}]{
                \texttt{\tiny{[GM$|$GM]}}
                \texttt{risti} \\
            }
        }
    }
    & & \\ \\

    \theutterance \stepcounter{utterance}  
    & & & \multicolumn{2}{p{0.3\linewidth}}{
        \cellcolor[rgb]{0.9,0.9,0.9}{
            \makecell[{{p{\linewidth}}}]{
                \texttt{\tiny{[GM$|$GM]}}
                {'word': 'syli', 'assignment': 'opponent'}
            }
        }
    }
    & & \\ \\

    \theutterance \stepcounter{utterance}  
    & & & \multicolumn{2}{p{0.3\linewidth}}{
        \cellcolor[rgb]{0.9,0.9,0.9}{
            \makecell[{{p{\linewidth}}}]{
                \texttt{\tiny{[GM$|$GM]}}
                {'hidden': {'team': ['klubi', 'minttu', 'koukku'], 'innocent': ['solmio', 'marmori', 'puku', 'punarinta', 'jää'], 'opponent': ['salakuljettaja'], 'assassin': ['kenguru']}, 'revealed': {'team': {'team': ['antarktis', 'oppilas', 'vallankumous', 'herätä', 'berliini', 'neliö'], 'innocent': ['kruunu', 'risti'], 'opponent': ['karhu', 'tiedemies'], 'assassin': []}, 'opponent': {'team': [], 'innocent': [], 'opponent': ['rivi', 'päivämäärä', 'elin', 'joukkovelkakirja', 'syli'], 'assassin': []}}}
            }
        }
    }
    & & \\ \\

    \theutterance \stepcounter{utterance}  
    & & \multicolumn{4}{p{0.6\linewidth}}{
        \cellcolor[rgb]{0.9,0.9,0.9}{
            \makecell[{{p{\linewidth}}}]{
                \texttt{\tiny{[P1$\langle$GM]}}
                \texttt{Sana neliö arvattiin oikein. Sana risti arvattiin, mutta se on viaton sana. Tiimikaverin vuoro loppui siellä.} \\
\texttt{Anna nyt toinen vihje, joka liittyy joihinkin jäljellä oleviin tiimisi sanoihin, ja luettelo niihin liittyvistä kohdesanoista. Muista aloittaa vihjeesi sanalla 'VIHJE: ', laita uusi rivi ja aloita pilkulla erotettu kohdesanojen luettelo sanalla 'KOHTEET: '. Huomaa: joitakin sanoja on poistettu luetteloista aiempiin pyyntöihin verrattuna.} \\
\\ 
\texttt{Jäljellä olevat tiimisi sanat ovat: klubi, minttu, koukku.} \\
\texttt{Vastustajan jäljellä olevat sanat ovat: syli, salakuljettaja.} \\
\texttt{Jäljellä olevat häiritsevät sanat ovat: solmio, marmori, puku, punarinta, jää.} \\
\texttt{Jäljellä olevat tapposanat ovat: kenguru.} \\
            }
        }
    }
    & \\ \\

    \theutterance \stepcounter{utterance}  
    & \multicolumn{4}{p{0.6\linewidth}}{
        \cellcolor[rgb]{0.9,0.9,0.9}{
            \makecell[{{p{\linewidth}}}]{
                \texttt{\tiny{[P1$\rangle$GM]}}
                \texttt{VIHJE: JÄRJESTÖ} \\
\texttt{KOHTEET: klubi, koukku} \\
            }
        }
    }
    & & \\ \\

    \theutterance \stepcounter{utterance}  
    & & & \multicolumn{2}{p{0.3\linewidth}}{
        \cellcolor[rgb]{0.9,0.9,0.9}{
            \makecell[{{p{\linewidth}}}]{
                \texttt{\tiny{[GM$|$GM]}}
                \texttt{järjestö} \\
            }
        }
    }
    & & \\ \\

    \theutterance \stepcounter{utterance}  
    & & & \multicolumn{2}{p{0.3\linewidth}}{
        \cellcolor[rgb]{0.9,0.9,0.9}{
            \makecell[{{p{\linewidth}}}]{
                \texttt{\tiny{[GM$|$GM]}}
                ['klubi', 'koukku']
            }
        }
    }
    & & \\ \\

    \theutterance \stepcounter{utterance}  
    & & & \multicolumn{2}{p{0.3\linewidth}}{
        \cellcolor[rgb]{0.9,0.9,0.9}{
            \makecell[{{p{\linewidth}}}]{
                \texttt{\tiny{[GM$|$GM]}}
                {'word': 'klubi', 'assignment': 'team'}
            }
        }
    }
    & & \\ \\

    \theutterance \stepcounter{utterance}  
    & & & \multicolumn{2}{p{0.3\linewidth}}{
        \cellcolor[rgb]{0.9,0.9,0.9}{
            \makecell[{{p{\linewidth}}}]{
                \texttt{\tiny{[GM$|$GM]}}
                {'word': 'koukku', 'assignment': 'team'}
            }
        }
    }
    & & \\ \\

    \theutterance \stepcounter{utterance}  
    & & \multicolumn{4}{p{0.6\linewidth}}{
        \cellcolor[rgb]{0.9,0.9,0.9}{
            \makecell[{{p{\linewidth}}}]{
                \texttt{\tiny{[GM$\rangle$P2]}}
                \texttt{Sana neliö arvattiin oikein. Sana risti arvattiin, mutta se on viaton sana. Sinun vuorosi loppui siellä.} \\
\texttt{Anna nyt toinen pilkulla erotettu luettelo vähintään 1:stä ja enintään 2 sanasta seuraavasta sanaluettelosta, jotka liittyvät parhaiten tai ovat läheisimmin yhteydessä sanaan 'järjestö'. Muista aloittaa vastauksesi ARVAUS:. Huomio: joitakin sanoja on poistettu luettelosta verrattuna aiempiin pyyntöihin.} \\
\\ 
\texttt{klubi, koukku, punarinta, solmio, salakuljettaja, jää, puku, kenguru, marmori, minttu} \\
            }
        }
    }
    & \\ \\

    \theutterance \stepcounter{utterance}  
    & & & \multicolumn{4}{p{0.6\linewidth}}{
        \cellcolor[rgb]{0.9,0.9,0.9}{
            \makecell[{{p{\linewidth}}}]{
                \texttt{\tiny{[GM$\langle$P2]}}
                \texttt{ARVAUS: klubi} \\
            }
        }
    }
     \\ \\

    \theutterance \stepcounter{utterance}  
    & & & \multicolumn{2}{p{0.3\linewidth}}{
        \cellcolor[rgb]{0.9,0.9,0.9}{
            \makecell[{{p{\linewidth}}}]{
                \texttt{\tiny{[GM$|$GM]}}
                ['klubi']
            }
        }
    }
    & & \\ \\

    \theutterance \stepcounter{utterance}  
    & & & \multicolumn{2}{p{0.3\linewidth}}{
        \cellcolor[rgb]{0.9,0.9,0.9}{
            \makecell[{{p{\linewidth}}}]{
                \texttt{\tiny{[GM$|$GM]}}
                {'word': 'klubi', 'assignment': 'team'}
            }
        }
    }
    & & \\ \\

    \theutterance \stepcounter{utterance}  
    & & & \multicolumn{2}{p{0.3\linewidth}}{
        \cellcolor[rgb]{0.9,0.9,0.9}{
            \makecell[{{p{\linewidth}}}]{
                \texttt{\tiny{[GM$|$GM]}}
                {'word': 'klubi', 'assignment': 'team'}
            }
        }
    }
    & & \\ \\

    \theutterance \stepcounter{utterance}  
    & & & \multicolumn{2}{p{0.3\linewidth}}{
        \cellcolor[rgb]{0.9,0.9,0.9}{
            \makecell[{{p{\linewidth}}}]{
                \texttt{\tiny{[GM$|$GM]}}
                {'word': 'salakuljettaja', 'assignment': 'opponent'}
            }
        }
    }
    & & \\ \\

    \theutterance \stepcounter{utterance}  
    & & & \multicolumn{2}{p{0.3\linewidth}}{
        \cellcolor[rgb]{0.9,0.9,0.9}{
            \makecell[{{p{\linewidth}}}]{
                \texttt{\tiny{[GM$|$GM]}}
                {'hidden': {'team': ['minttu', 'koukku'], 'innocent': ['solmio', 'marmori', 'puku', 'punarinta', 'jää'], 'opponent': [], 'assassin': ['kenguru']}, 'revealed': {'team': {'team': ['antarktis', 'oppilas', 'vallankumous', 'herätä', 'berliini', 'neliö', 'klubi'], 'innocent': ['kruunu', 'risti'], 'opponent': ['karhu', 'tiedemies'], 'assassin': []}, 'opponent': {'team': [], 'innocent': [], 'opponent': ['rivi', 'päivämäärä', 'elin', 'joukkovelkakirja', 'syli', 'salakuljettaja'], 'assassin': []}}}
            }
        }
    }
    & & \\ \\

    \theutterance \stepcounter{utterance}  
    & & \multicolumn{4}{p{0.6\linewidth}}{
        \cellcolor[rgb]{0.9,0.9,0.9}{
            \makecell[{{p{\linewidth}}}]{
                \texttt{\tiny{[P1$\langle$GM]}}
                \texttt{Sana klubi arvattiin oikein. Tiimikaverin vuoro loppui siellä.} \\
\texttt{Anna nyt toinen vihje, joka liittyy joihinkin jäljellä oleviin tiimisi sanoihin, ja luettelo niihin liittyvistä kohdesanoista. Muista aloittaa vihjeesi sanalla 'VIHJE: ', laita uusi rivi ja aloita pilkulla erotettu kohdesanojen luettelo sanalla 'KOHTEET: '. Huomaa: joitakin sanoja on poistettu luetteloista aiempiin pyyntöihin verrattuna.} \\
\\ 
\texttt{Jäljellä olevat tiimisi sanat ovat: minttu, koukku.} \\
\texttt{Vastustajan jäljellä olevat sanat ovat: salakuljettaja.} \\
\texttt{Jäljellä olevat häiritsevät sanat ovat: solmio, marmori, puku, punarinta, jää.} \\
\texttt{Jäljellä olevat tapposanat ovat: kenguru.} \\
            }
        }
    }
    & \\ \\

    \theutterance \stepcounter{utterance}  
    & \multicolumn{4}{p{0.6\linewidth}}{
        \cellcolor[rgb]{0.9,0.9,0.9}{
            \makecell[{{p{\linewidth}}}]{
                \texttt{\tiny{[P1$\rangle$GM]}}
                \texttt{VIHJE: MAKU} \\
\texttt{KOHTEET: minttu} \\
            }
        }
    }
    & & \\ \\

    \theutterance \stepcounter{utterance}  
    & & & \multicolumn{2}{p{0.3\linewidth}}{
        \cellcolor[rgb]{0.9,0.9,0.9}{
            \makecell[{{p{\linewidth}}}]{
                \texttt{\tiny{[GM$|$GM]}}
                \texttt{maku} \\
            }
        }
    }
    & & \\ \\

    \theutterance \stepcounter{utterance}  
    & & & \multicolumn{2}{p{0.3\linewidth}}{
        \cellcolor[rgb]{0.9,0.9,0.9}{
            \makecell[{{p{\linewidth}}}]{
                \texttt{\tiny{[GM$|$GM]}}
                ['minttu']
            }
        }
    }
    & & \\ \\

    \theutterance \stepcounter{utterance}  
    & & & \multicolumn{2}{p{0.3\linewidth}}{
        \cellcolor[rgb]{0.9,0.9,0.9}{
            \makecell[{{p{\linewidth}}}]{
                \texttt{\tiny{[GM$|$GM]}}
                {'word': 'minttu', 'assignment': 'team'}
            }
        }
    }
    & & \\ \\

    \theutterance \stepcounter{utterance}  
    & & & \multicolumn{2}{p{0.3\linewidth}}{
        \cellcolor[rgb]{0.9,0.9,0.9}{
            \makecell[{{p{\linewidth}}}]{
                \texttt{\tiny{[GM$|$GM]}}
                \texttt{opponent has won} \\
            }
        }
    }
    & & \\ \\

\end{supertabular}
}

\end{document}
